\documentclass[11pt,oneside]{amsart}
\usepackage[T1]{fontenc}
\usepackage{lmodern}
\usepackage[hmargin=1.5in,vmargin=1.25in]{geometry}
\linespread{1.12}
\setlength{\emergencystretch}{1.5em}
\usepackage{etoolbox}
\newlength{\abstractmargin}
\setlength{\abstractmargin}{6pc}
\makeatletter
\patchcmd{\abstract}{\leftmargin3pc}{\leftmargin\abstractmargin}{}{\PackageError{paper}{abstract patch failed}{}}
\makeatother
\usepackage{amssymb,amsmath,mathtools}
\usepackage[expansion=false]{microtype}
\usepackage{needspace}
\usepackage{xcolor}
\definecolor{linkblue}{RGB}{0,58,112}
\usepackage[colorlinks=true,allcolors=linkblue]{hyperref}
\hypersetup{
  pdftitle={Property F implies weak integrality for pseudo-unitary braided fusion categories},
  pdfauthor={},
  pdfsubject={Braided fusion categories and finite braid group images},
  pdfkeywords={property F, weak integrality, fusion categories, Galois symmetry}
}

\newtheorem{theorem}{Theorem}[section]
\newtheorem{lemma}[theorem]{Lemma}
\newtheorem{proposition}[theorem]{Proposition}
\newtheorem{corollary}[theorem]{Corollary}
\newtheorem*{conjecture}{Conjecture}
\theoremstyle{remark}
\newtheorem{remark}[theorem]{Remark}

\newcommand{\C}{\mathcal C}
\newcommand{\D}{\mathcal D}
\newcommand{\M}{\mathcal M}
\newcommand{\one}{\mathbf 1}
\newcommand{\FP}{\operatorname{FPdim}}
\newcommand{\Tr}{\operatorname{Tr}}
\newcommand{\tr}{\operatorname{tr}}
\newcommand{\ptr}{\operatorname{ptr}}
\newcommand{\Hom}{\operatorname{Hom}}
\newcommand{\Hopf}{\operatorname{Hopf}}
\newcommand{\Rep}{\operatorname{Rep}}
\newcommand{\Vect}{\mathrm{Vec}}
\newcommand{\sVect}{\mathrm{sVec}}
\newcommand{\Z}{\mathbb Z}
\newcommand{\hs}{\hat\sigma}
\newcommand{\htau}{\hat\tau}
\newcommand{\sourcecite}[3]{\nocite{#1}\href{#3}{[#1, #2]}}

\title[Property F and weak integrality]{Property F implies weak integrality\\
for pseudo-unitary braided fusion categories}
\date{October 15, 2026}

\begin{document}
\pagestyle{plain}

\begin{abstract}
We prove that every pseudo-unitary braided fusion category with property F is
weakly integral. Here property F means that all braid group representations
arising from the braiding have finite image. More generally, the conclusion holds
for every braided fusion category with property F that has a pseudo-unitary Galois
conjugate. This establishes one direction of the Naidu--Rowell property F
conjecture for all pseudo-unitary, and hence all unitary, braided fusion
categories.

The proof is by contradiction. Using Galois symmetry of the modular data, we
construct a non-invertible, self-dual simple object $Y$ in a pseudo-unitary
Drinfeld centre, lying in a subcategory with property F. A multi-strand trace
inequality, combined with the balancing axiom, forces every simple constituent of
$Y\otimes Y$ to have twist $1$ or $-1$. A three-strand Markov identity, together
with Galois symmetry, then forces $Y$ to be invertible, giving a contradiction.
\end{abstract}

\maketitle

\section{Introduction}

Let $\C$ be a braided fusion category over $\mathbb C$. For an object $X$ and an
integer $n\ge2$, the braiding defines a representation
\[
\rho_{X,n}:B_n\longrightarrow\operatorname{Aut}(X^{\otimes n}),\qquad
\sigma_i\longmapsto
\mathrm{id}_X^{\otimes(i-1)}\otimes c_{X,X}\otimes
\mathrm{id}_X^{\otimes(n-i-1)}.
\]
We write $G_{X,n}$ for its image. The category $\C$ has \emph{property F} if
$G_{X,n}$ is finite for every $X$ and $n$. It is \emph{weakly integral} if
$\FP(\C)\in\Z$.

\begin{conjecture}[Naidu--Rowell {\cite{NR}}]
A braided fusion category has property F if and only if it is weakly integral.
\end{conjecture}

Thus the conjecture predicts that a single arithmetic condition on the category
determines finiteness of its braid group images, for every object and every
number of strands.

The implication from weak integrality to property F is known for weakly
group-theoretical braided fusion categories \cite[Theorem~7.2]{GN}, extending the
group-theoretical case \cite{ERW}. Etingof, Nikshych and Ostrik asked whether there
exists a weakly integral fusion category that is not weakly group-theoretical
\sourcecite{ENO2}{Question~2}{https://arxiv.org/html/0809.3031\#S11.p3.1.1}; a
negative answer would establish that implication in general. In the other direction, infinite braid group images have
been established for several families of quantum-group categories
\cite{FLW,LRW}. We prove the following general converse under a
pseudo-unitarity hypothesis. Recall that a fusion category is
\emph{pseudo-unitary} if its global dimension $\dim(\C)$ equals its
Frobenius--Perron dimension $\FP(\C)$; unitary fusion categories are
pseudo-unitary.

\begin{theorem}\label{thm:main}
Let $\C$ be a braided fusion category with a pseudo-unitary Galois conjugate.
If $\C$ has property F, then $\C$ is weakly integral. In particular, every
pseudo-unitary braided fusion category with property F is weakly integral.
\end{theorem}

In particular, the theorem applies to all unitary braided fusion categories. Equivalently, a non-weakly-integral category
in this class admits a braid group representation with infinite image. In fact
the result holds object by object: every simple object $X$ with
$\FP(X)^2\notin\Z$ in such a category already has infinite braid group images
(Corollary~\ref{cor:object}), and the modular data determine an explicit number
of strands at which infinitude is certified (Corollary~\ref{cor:strands}).
The trace inequality behind this (Theorem~\ref{thm:ms}) needs no pseudo-unitarity,
and Section~\ref{sec:conseq} applies it outside the hypothesis of Theorem~\ref{thm:main}.

The proof is by contradiction. Starting from a category with property F that is
not weakly integral, we construct a pseudo-unitary modular category $\M$ and a
self-dual, non-invertible simple object $Y$ in $\M$. The object $Y$ lies in a
fusion subcategory with property F, and it is the image of the unit label under
the Galois symmetry of the modular data.

The two main arguments give incompatible conclusions about $Y$. First,
finiteness of its braid group images bounds the sum of the absolute values of
certain categorical traces in a Galois-conjugate category. The balancing axiom
forces equality in this bound, which implies that every simple constituent of
$Y\otimes Y$ has twist $1$ or $-1$. Second, a three-strand trace identity,
together with Galois symmetry, shows that an object obtained from the unit label
cannot have this twist pattern unless it is invertible. This contradicts the
choice of $Y$.

Section~\ref{sec:prelim} records the transfer properties and Galois identities
used in the proof. Section~\ref{sec:red} reduces the theorem to a pseudo-unitary
Drinfeld centre and locates $Y$ in a subcategory with property F.
Sections~\ref{sec:ms} and~\ref{sec:align} establish the trace inequality and its
twist consequence. Section~\ref{sec:QO} proves the Markov identity and the
obstruction for non-invertible images of the unit label. The proof of
Theorem~\ref{thm:main} is completed in Section~\ref{sec:main}.

\section{Preliminaries}\label{sec:prelim}

In the sequel, all fusion categories are over $\mathbb C$, and we suppress
associativity and unit constraints. We write $\one$ for the tensor unit and
$N_{XY}^Z$ for the multiplicity of a simple object $Z$ in $X\otimes Y$. A
\emph{simple constituent} of a tensor product is a simple object occurring in
it (physicists call these \emph{channels}).

\subsection{Transfer of property F}

\begin{lemma}[Direct summands]\label{lem:summands}
Let $Y$ be a direct summand of an object $Z$ in a braided fusion category.
Then $G_{Y,n}$ is a quotient of $G_{Z,n}$ for every $n\ge2$.
\end{lemma}

\begin{proof}
Choose morphisms $\iota:Y\to Z$ and $\pi:Z\to Y$ with
$\pi\iota=\mathrm{id}_Y$. Naturality of the braiding gives
\[
\rho_{Y,n}(b)=\pi^{\otimes n}\rho_{Z,n}(b)\iota^{\otimes n}
\qquad(b\in B_n).
\]
Thus $\ker\rho_{Z,n}\subseteq\ker\rho_{Y,n}$, which gives the asserted quotient.
\end{proof}

\begin{lemma}[Braided functors and subcategories]\label{lem:functors}
Property F passes to fusion subcategories and along dominant braided tensor
functors.
\end{lemma}

\begin{proof}
The subcategory assertion is immediate. If $F:\C\to\D$ is a braided tensor
functor, its tensor structure identifies $\rho_{F(X),n}(b)$ with
$F(\rho_{X,n}(b))$. Hence $G_{F(X),n}$ is a quotient of $G_{X,n}$. If $F$ is
dominant, every object of $\D$ is a direct summand of some $F(X)$, so
Lemma~\ref{lem:summands} applies.
\end{proof}

\begin{lemma}[Reverse braidings and tensor products]\label{lem:products}
Let $\C$ and $\D$ be braided fusion categories. The category $\C$ has property F
if and only if $\C^{\rm rev}$ does. Moreover, $\C\boxtimes\D$ has property F
if and only if both $\C$ and $\D$ do.
\end{lemma}

\begin{proof}
The reverse braiding satisfies $c^{\rm rev}_{X,X}=c_{X,X}^{-1}$, and
$\sigma_i\mapsto\sigma_i^{-1}$ defines an automorphism of $B_n$. For an external
tensor product $X\boxtimes Y$, the braid image is a quotient of a subgroup of
$G_{X,n}\times G_{Y,n}$. Every object
$\bigoplus_iX_i\boxtimes Y_i$ is a direct summand of
$(\bigoplus_iX_i)\boxtimes(\bigoplus_iY_i)$, so Lemma~\ref{lem:summands}
proves the forward implication. The converse follows from the braided
embeddings of the two factors.
\end{proof}

\begin{lemma}[Closure]\label{lem:closure}
Let $X$ be an object of a braided fusion category, and let $\langle X\rangle$ be
the fusion subcategory it generates. If $G_{X,n}$ is finite for every $n$, then
$\langle X\rangle$ has property F.
\end{lemma}

\begin{proof}
First let $Z=m\one\oplus X$ with $m\ge1$. The braid group permutes the direct-sum
components of $Z^{\otimes n}$ indexed by sequences of summand labels, and the
pure braid group, which has finite index, preserves each of them. Since the
braiding with $\one$ is trivial, a pure braid acts on the component with $k$
copies of $X$ through the braid obtained by deleting the unit strands, so through
the finite group $G_{X,k}$. Hence $G_{Z,n}$ is finite. Next, $G_{Z^{\otimes r},n}$
is the image of a subgroup of $B_{rn}$ (cabled braids) acting on
$Z^{\otimes rn}$, so it is finite too. Finally, every object of $\langle X\rangle$
is a direct summand of some $X^{\otimes k}$ (the dual $X^*$ occurs in a tensor
power of $X$), hence of some $(m\one\oplus X)^{\otimes r}$. Apply
Lemma~\ref{lem:summands}.
\end{proof}

\begin{lemma}[Galois invariance]\label{lem:galois-transfer}
Property F and weak integrality are invariant under Galois conjugation.
\end{lemma}

\begin{proof}
By Ocneanu rigidity, a fusion category admits a model over a number field
\cite[Theorem~2.30]{ENO}, and it has only finitely many braidings
\cite[Remark~2.33]{ENO}. The braidings are therefore the solutions of a
zero-dimensional system of polynomial equations over that field, so they are
defined over a finite extension of it. Galois conjugation acts entrywise on braid matrices and preserves
equality and distinctness of matrices. It therefore induces an isomorphism
between the original and conjugate braid images. Weak integrality depends
only on the fusion ring, which Galois conjugation preserves.
\end{proof}

\subsection{Dimensions and the centre of a braided category}

For a spherical fusion category, $d_X$ denotes the categorical dimension of
$X$ and $\Tr$ the categorical trace. Simple categorical dimensions and their
Galois conjugates are real and nonzero
\sourcecite{ENO}{Theorem~2.3, Remark~2.5 and Corollary~2.10}{https://annals.math.princeton.edu/wp-content/uploads/annals-v162-n2-p01.pdf\#page=5}:
the squared norm $|X|^2$ of a simple object is totally positive, and in a
spherical category $|X|^2=d_X^2$. The global dimension is $\dim(\C)=\sum_Xd_X^2$, with $X$ ranging over the
simple objects. A pseudo-unitary fusion category has a unique spherical
structure with $d_X=\FP(X)$ for every simple $X$
\sourcecite{ENO}{Proposition~8.23}{https://annals.math.princeton.edu/wp-content/uploads/annals-v162-n2-p01.pdf\#page=42};
we always use this structure on a pseudo-unitary category. A \emph{modular
category} is a non-degenerate braided fusion category with a spherical
structure and its associated ribbon twist.

Let $\D$ be a braided fusion category. Its \emph{M\"uger centre} $\D'$ is the
fusion subcategory of objects $X$ with $c_{Y,X}c_{X,Y}=\mathrm{id}$ for every
object $Y$. There are fully faithful braided tensor
functors
\[
\begin{aligned}
\iota_+:\D&\longrightarrow\mathcal Z(\D),
&X&\longmapsto(X,c_{X,-}),\\
\iota_-:\D^{\rm rev}&\longrightarrow\mathcal Z(\D),
&X&\longmapsto(X,c^{-1}_{-,X}).
\end{aligned}
\]
Their images, denoted $\D_+$ and $\D_-$, centralise one another. The Drinfeld
centre $\mathcal Z(\D)$ is non-degenerate and satisfies
$\FP(\mathcal Z(\D))=\FP(\D)^2$. It is modular if $\D$ is spherical
\cite[Theorem~7.16.6 and Corollaries~8.20.13 and~8.20.14]{EGNO}. It is
pseudo-unitary if $\D$ is: indeed $\dim(\mathcal Z(\D))=\dim(\D)^2$
\cite[Proposition~9.3.4]{EGNO}, so $\dim(\D)=\FP(\D)$ implies
$\dim(\mathcal Z(\D))=\FP(\mathcal Z(\D))$. We write
$\mathcal A\vee\mathcal B$ for the fusion subcategory generated by two fusion
subcategories $\mathcal A$ and $\mathcal B$.

\begin{lemma}\label{lem:centre}
Inside $\mathcal Z(\D)$, the subcategory $\D_+\vee\D_-$ is the centraliser of
$\iota_+(\D')$. If $\D$ has property F, then $\D_+\vee\D_-$ has property F.
\end{lemma}

\begin{proof}
The two images centralise one another, so $\D_-\subseteq\D_+'$, where the prime
now denotes the centraliser in $\mathcal Z(\D)$. The dimension formula for
centralisers gives
\[
\FP(\D_+')=\frac{\FP(\mathcal Z(\D))}{\FP(\D_+)}
=\FP(\D)=\FP(\D_-)
\]
\sourcecite{DGNO}{Theorem~3.14}{https://arxiv.org/html/0906.0620\#S3.Thmtheorem14}. Thus $\D_+'=\D_-$, and
$\D_+\cap\D_-=\iota_+(\D')$. Since $\D_-'=\D_+''=\D_+$, this gives
\[
(\D_+\vee\D_-)'=\D_+'\cap\D_-'=\D_-\cap\D_+=\iota_+(\D').
\]
Taking centralisers, and using the double-centraliser theorem
\sourcecite{DGNO}{Theorem~3.10(ii)}{https://arxiv.org/html/0906.0620\#S3.Thmtheorem10}
in the non-degenerate category $\mathcal Z(\D)$, proves the first assertion.

Since the two images centralise one another, the tensor product functor
\[
\D\boxtimes\D^{\rm rev}\longrightarrow\D_+\vee\D_-,\qquad
X\boxtimes Y\longmapsto\iota_+(X)\otimes\iota_-(Y),
\]
is braided and dominant. Lemmas~\ref{lem:functors} and~\ref{lem:products}
prove the second assertion.
\end{proof}

\subsection{Galois symmetry of modular data}

Let $\M$ be modular. We normalise the twist by
\[
\ptr_2(c_{A,A})=\theta_A\,\mathrm{id}_A
\]
for a simple object $A$, where $\ptr_2$ is the right partial trace. Twists are
roots of unity \cite[Theorem~5.1]{E}, so every Galois conjugate of a twist has absolute value
one. For simple objects $X$ and $A$, set
\begin{equation}\label{eq:hopf}
\Hopf(X,A):=\Tr_{X\otimes A}(c_{A,X}c_{X,A})
=\sum_ZN_{XA}^Z\,d_Z\,\frac{\theta_Z}{\theta_X\theta_A}.
\end{equation}
The second equality is the balancing axiom. We normalise the $S$-matrix by
\[
s_{XY}=\frac{\Hopf(X,Y)}{\sqrt{\dim(\M)}},\qquad
\chi_Y(X)=\frac{s_{XY}}{s_{\one Y}},
\]
using the positive square root of the global dimension. The characters of the
fusion ring $K_0(\M)$ are precisely the $\chi_Y$. Fix a cyclotomic field $K$
containing the entries of $s$, the twists, and $i$
\cite[Theorem~7.1]{NS}, \cite[Theorem~II]{DLN}; the positive square root of
$\dim(\M)$ is an admissible normalisation there.
Galois automorphisms are taken in
$\operatorname{Gal}(\overline{\mathbb Q}/\mathbb Q)$; their action on modular
data depends only on their restriction to $K$.

\Needspace{6\baselineskip}
For a Galois automorphism $\tau$, the character $\tau\circ\chi_Y$ is again a
fusion-ring character. Thus there is a permutation $\htau$ of the simple
objects such that $\tau\circ\chi_Y=\chi_{\htau(Y)}$.

The permutation $\htau$ acts on simple-object labels through their fusion-ring
characters. In particular, $\htau(\one)$ need not be the tensor unit. This
action on labels is distinct from passing to the Galois-conjugate category
$\M^\tau$, whose tensor unit is still $\one$. Both actions appear below: the
conjugate category in Section~\ref{sec:ms}, and the permutation of labels
everywhere else.

\begin{proposition}[Galois symmetry {\cite{CG}, \cite[\S4.1]{DLN}}]\label{prop:galois}
Let $\M$ be a modular category and $\tau$ a Galois automorphism. For every
simple object $X$,
\begin{align}
\tau(d_X)&=\chi_{\htau(\one)}(X),\label{eq:galois-character}\\
\tau(d_X)&=\pm\frac{d_{\htau(X)}}{d_{\htau(\one)}},\label{eq:galois-dimension}\\
\tau(\dim(\M))&=\frac{\dim(\M)}{d_{\htau(\one)}^2}.\label{eq:galois-global}
\end{align}
Moreover, $\htau(\one)$ is self-dual, and
\begin{equation}\label{eq:galois-inverse}
\widehat{\tau^{-1}}=\htau^{-1}.
\end{equation}
\end{proposition}

\begin{proof}
Equation~\eqref{eq:galois-character} follows from $\chi_{\one}(X)=d_X$.
The signed-permutation formula \cite[(4.3)]{DLN}, see also
\cite[Theorem~10.1]{ENO},
\[
\tau(s_{XY})=\varepsilon_\tau(X)s_{\htau(X),Y},\qquad
\varepsilon_\tau(X)\in\{1,-1\},
\]
together with $d_X=s_{X\one}/s_{\one\one}$ and
$\dim(\M)=s_{\one\one}^{-2}$, gives
\eqref{eq:galois-dimension} and~\eqref{eq:galois-global}.
For self-duality, use
$\chi_{Y^*}=\overline{\chi_Y}$ and the fact that complex conjugation commutes
with $\tau$ on $K$. Since $\chi_{\one}$ is real-valued, so is
$\tau\circ\chi_{\one}$, which implies that $\htau(\one)$ is self-dual.
The inverse-permutation assertion follows by composing the actions of
$\tau$ and $\tau^{-1}$ on characters.
\end{proof}

\begin{lemma}[Galois column]\label{lem:galcol}
If $Y=\htau(\one)$, then
$\Hopf(X,Y)=d_Y\,\tau(d_X)$ for every simple $X$.
\end{lemma}

\begin{proof}
By the chosen normalisation and \eqref{eq:galois-character},
\[
\Hopf(X,Y)=d_Y\frac{s_{XY}}{s_{\one Y}}
=d_Y\chi_Y(X)=d_Y\tau(d_X).\qedhere
\]
\end{proof}

\section{Reduction to a pseudo-unitary Drinfeld centre}\label{sec:red}

We now pass to a modular category, so that the Galois identities of
Section~\ref{sec:prelim} are available. The reduction also identifies a
subcategory with property F inside this modular category, and places the simple
object needed for the trace arguments in that subcategory. The roles are as
follows:
\begin{itemize}
\item $\M=\mathcal Z(\C)$ provides modular data and Galois symmetry;
\item $\mathcal B:=\C_+\vee\C_-\subseteq\M$ provides property F;
\item $Y\in\mathcal B$ is the object to which both later arguments apply.
\end{itemize}
The argument never needs property F for the whole Drinfeld centre.

\begin{proposition}\label{prop:red}
Let $\C$ be a pseudo-unitary braided fusion category that has property F and is
not weakly integral. Then $\M:=\mathcal Z(\C)$ is a pseudo-unitary modular
category that is not weakly integral, and $\mathcal B=\C_+\vee\C_-$ has
property F.
\end{proposition}

\begin{proof}
Lemma~\ref{lem:centre} shows that $\mathcal B$ has property F, and $\M$ is
pseudo-unitary and modular. Recall that $\FP(\M)=\FP(\C)^2$. Since $\C$ is
pseudo-unitary, $\FP(\C)=\dim(\C)$. This number is a totally positive algebraic
integer: each of its Galois conjugates is the positive global dimension of a
Galois-conjugate fusion category
\sourcecite{ENO}{Theorem~2.3 and Remark~2.5}{https://annals.math.princeton.edu/wp-content/uploads/annals-v162-n2-p01.pdf\#page=5}.
If $\FP(\C)^2$ were an integer $m$, then an irrational $\FP(\C)=\sqrt m$ would
have the negative conjugate $-\sqrt m$, contradicting total positivity. Thus
$\FP(\C)$ would be rational, hence an integer, contrary to the hypothesis.
\end{proof}

The next lemma locates every Galois image of the unit label. It says that such
an object braids trivially with every simple object of integer dimension.

\begin{lemma}[Integral centraliser]\label{lem:intcent}
Let $\M$ be a pseudo-unitary modular category, $\tau$ a Galois automorphism,
and $Y=\htau(\one)$. Then $Y$ centralises every simple object $z$ of $\M$ with
$d_z\in\Z$. In particular, $Y$ centralises every integral fusion subcategory
of $\M$.
\end{lemma}

\begin{proof}
Since $d_z\in\Z$, Lemma~\ref{lem:galcol} gives
$\Hopf(z,Y)=d_Y\tau(d_z)=d_Yd_z$. On the other hand, by~\eqref{eq:hopf},
\[
\Hopf(z,Y)=\sum_WN_{zY}^W\,d_W\,m_W,\qquad
m_W:=\frac{\theta_W}{\theta_z\theta_Y},
\]
where each $m_W$ has absolute value one, each coefficient $N_{zY}^Wd_W$ is
positive by pseudo-unitarity, and the coefficients sum to $d_zd_Y$. Equality
with $d_zd_Y$ therefore forces $m_W=1$ for every $W$. The monodromy
$c_{Y,z}c_{z,Y}$ acts on the $W$-isotypic part of $z\otimes Y$ by $m_W$, so it
is the identity.
\end{proof}

Now let $\C$ be as in Proposition~\ref{prop:red}, and let $\M$ and $\mathcal B$
be as there. Since the algebraic integer $\dim(\M)=\FP(\M)$ is not an integer,
it is irrational. Choose a Galois automorphism $\tau$ with
$\tau(\dim(\M))\ne\dim(\M)$, and set
\[
Y:=\htau(\one),\qquad \sigma:=\tau^{-1}.
\]
By~\eqref{eq:galois-global}, $d_Y^2=\dim(\M)/\tau(\dim(\M))\ne1$, so $Y$ is
non-invertible. By Proposition~\ref{prop:galois}, $Y$ is self-dual.

\begin{lemma}\label{lem:Yloc}
The object $Y$ belongs to $\mathcal B$.
\end{lemma}

\begin{proof}
By Lemma~\ref{lem:centre}, $\mathcal B$ is the centraliser of $\iota_+(\C')$.
The M\"uger centre $\C'$ is a symmetric fusion category. By Deligne's
theorem \cite[Corollaire~0.8]{Del} it is equivalent to $\Rep(G,z)$ for a finite group $G$ and a
central element $z$ of order at most two, so the Frobenius--Perron dimension of
each simple object is the dimension of a representation of $G$. Thus $\C'$ is
integral.
Since $\iota_+$ is a tensor functor, it preserves Frobenius--Perron dimensions,
and these agree with the categorical dimensions in the pseudo-unitary category
$\M$. So $\iota_+(\C')$ is an integral fusion subcategory of $\M$, and
Lemma~\ref{lem:intcent} applies.
\end{proof}

We have therefore obtained a non-invertible, self-dual simple object $Y$ in a
pseudo-unitary modular category, contained in a subcategory with property F.
The next two sections use this last property to constrain the twists of the
simple constituents of $Y\otimes Y$.

\section{A multi-strand trace inequality}\label{sec:ms}

Let $\mathcal B$ be a spherical braided fusion category, and let
$Y\in\mathcal B$ be simple; self-duality is not needed in this section. Fix a
Galois automorphism $\sigma$. The conjugate category $\mathcal B^\sigma$ has the
same fusion rules, dimensions $\sigma(d_W)$ and twists $\sigma(\theta_W)$, and by
Lemma~\ref{lem:galois-transfer} its braid images are isomorphic to those of
$\mathcal B$. Throughout this section, categorical traces $\Tr$ are taken in
$\mathcal B^\sigma$, and $\tr$ denotes the ordinary trace on a
finite-dimensional vector space. Put
\[
q:=\sigma(d_Y)\in\mathbb R\setminus\{0\},\qquad t:=|q|.
\]

In the Galois-conjugate category, categorical dimensions can have either sign.
Consequently the traces $u_\lambda$ of the monodromy eigenspaces defined below
may cancel substantially: their sum is $q^2$, whereas the sum
$V=\sum_\lambda|u_\lambda|$ of their absolute values can be much larger. We bound
$V$ using finiteness of the braid images, or more precisely a unitarity
property that finiteness implies.

Partition the simple constituents of $Y\otimes Y$ into classes according to their twists
in $\mathcal B$. For a class $\lambda$, let $\theta_\lambda$ be its common
twist. In $\mathcal B^\sigma$, the monodromy $c_{Y,Y}^2$ acts on this class by
$\sigma(\theta_\lambda/\theta_Y^2)$. Let $E_\lambda$ be the corresponding
spectral projection, and define
\[
u_\lambda:=\Tr(E_\lambda)
=\sum_{c\in\lambda}N_{YY}^c\,\sigma(d_c),\qquad
S_\pm:=\sum_{\pm u_\lambda>0}|u_\lambda|,\qquad
V:=\sum_\lambda|u_\lambda|=S_++S_-.
\]
These numbers are real, and $S_+-S_-=q^2$. Set
$E=\sum_{u_\lambda>0}E_\lambda$. On $Y^{\otimes n}$, let $P_i$ be $E$ acting
on strands $i,i+1$, and put
\[
Q_i^+=P_i,\qquad Q_i^-=1-P_i,\qquad
w^+=S_+,\qquad w^-=-S_-.
\]
Thus $w^\pm$ is the categorical trace of $Q_i^\pm$ on the two strands where
it acts. Each $P_i$ is a polynomial in $\rho_{Y,n}(\sigma_i^2)$.

\begin{lemma}[Trace factorisation]\label{lem:fact}
For signs $\varepsilon_1,\ldots,\varepsilon_{n-1}$, let $F$ be a product in
which each factor $Q_i^{\varepsilon_i}$ occurs exactly once, in any order. Then
\[
\Tr(F)=\frac{\prod_{i=1}^{n-1}w^{\varepsilon_i}}{q^{n-2}}.
\]
\end{lemma}

\begin{proof}
We induct on $n$, the case $n=2$ being the definition of $w^\pm$. The only
factor meeting strand $n$ is $Q_{n-1}^{\varepsilon_{n-1}}$. Write
$F=F_1Q_{n-1}^{\varepsilon_{n-1}}F_2$, where $F_1$ and $F_2$ act on the first
$n-1$ strands. Simplicity of $Y$ and the definition of $w^\pm$ imply
\[
\ptr_n(Q_{n-1}^{\pm})=\frac{w^\pm}{q}\,
\mathrm{id}_{Y^{\otimes(n-1)}}.
\]
The bimodule property of the partial trace therefore gives
\[
\Tr(F)=\frac{w^{\varepsilon_{n-1}}}{q}\,\Tr(F_1F_2).
\]
Apply the induction hypothesis to $F_1F_2$.
\end{proof}

\Needspace{7\baselineskip}
\begin{theorem}\label{thm:ms}
Fix $n\ge2$, and set $M_n:=\sum_WN_{Y^{\otimes n}}^W\,|\sigma(d_W)|$. Suppose that
each multiplicity space $\Hom(W,Y^{\otimes n})$ in $\mathcal B^\sigma$ carries a
positive definite inner product for which the adjacent monodromies
$\rho^\sigma_{Y,n}(\sigma_i^2)$, $1\le i\le n-1$, are all unitary. This holds,
for example, if $G_{Y,n}$ is finite. Then
\begin{equation}\label{eq:ms}
V^{n-1}\le t^{n-2}M_n.
\end{equation}
\end{theorem}

\begin{proof}
Fix $n$, and write
\[
N_W=N^W_{Y^{\otimes n}},\qquad H_W=\Hom(W,Y^{\otimes n}).
\]
If $G_{Y,n}$ is finite, then so is the braid image in $\mathcal B^\sigma$, and
averaging over it gives an inner product on each $H_W$ that is invariant under
the entire braid image; in particular the hypothesis holds. In general, each
$P_i$ is a spectral projection of the unitary monodromy operator
$\rho^\sigma_{Y,n}(\sigma_i^2)$, so every $P_i$ is an orthogonal projection.

The idea is to follow one unitary operator through two trace calculations. Set
\[
J_i:=Q_i^+-Q_i^-=2P_i-1,\qquad U:=J_1J_2\cdots J_{n-1}.
\]
The involution $J_i$ reverses the sign of the negative class traces, so its
categorical trace on the two strands where it acts is $w^+-w^-=S_++S_-=V$.
Averaging gives positive inner products on the multiplicity spaces, while the
categorical trace weights the resulting ordinary traces by the possibly negative
dimensions $\sigma(d_W)$.

Each $J_i$ is a unitary involution, so $U$ is unitary on every $H_W$; the $J_i$
need not commute for this. Writing $U_W$ for the action of $U$ on $H_W$, we have
$|\tr(U_W)|\le N_W$, and hence
\begin{equation}\label{eq:trace-bound}
|\Tr(U)|=\Bigl|\sum_W\sigma(d_W)\tr(U_W)\Bigr|
\le\sum_W|\sigma(d_W)|\,N_W=M_n.
\end{equation}

On the other hand, expanding each $J_i=Q_i^+-Q_i^-$ writes $U$ as a signed sum
of products in which each $Q_i^{\pm}$ occurs exactly once. Lemma~\ref{lem:fact}
applies to every term, and gives
\[
\Tr(U)=\frac{(w^+-w^-)^{n-1}}{q^{n-2}}=\frac{V^{n-1}}{q^{n-2}}.
\]
Taking absolute values and comparing with \eqref{eq:trace-bound} proves
\eqref{eq:ms}.
\end{proof}

The theorem becomes effective once $M_n$ is bounded. For every simple $W$,
$\sigma\circ d$ is a character of the fusion ring, so $\sigma(d_W)$ is an
eigenvalue of the fusion matrix of $W$. By Perron--Frobenius theory its absolute
value is at most $\FP(W)$. Writing $f=\FP(Y)$, this gives
\begin{equation}\label{eq:Mn}
M_n\le\sum_WN^W_{Y^{\otimes n}}\FP(W)=f^n.
\end{equation}

\begin{corollary}[Strand bound]\label{cor:strands}
If $\mathcal B$ has property F, then $V\le t\,\FP(Y)$. Conversely, suppose
that $r:=V/(t\,\FP(Y))>1$. Then $G_{Y,n}$ is infinite for every
\[
n\ge n_0:=\left\lfloor\frac{\log(\FP(Y)/t)}{\log r}\right\rfloor+2.
\]
\end{corollary}

\begin{proof}
If $G_{Y,n}$ is finite, then \eqref{eq:ms} and \eqref{eq:Mn} give
$V^{n-1}\le t^{n-2}f^n$, that is, $r^{n-1}\le f/t$. Under property F this holds
for all $n$, and letting $n\to\infty$ gives $r\le1$. If $r>1$, it fails as soon
as $n-1>\log(f/t)/\log r$, which holds for every $n\ge n_0$.
\end{proof}

\section{Twists of the simple constituents}\label{sec:align}

Return to the categories $\C$, $\M$ and $\mathcal B$, the automorphisms $\sigma$ and
$\tau=\sigma^{-1}$, and the object $Y=\htau(\one)$ constructed in
Section~\ref{sec:red}. Use the twist classes and weights $u_\lambda$ of
Section~\ref{sec:ms}, with $\mathcal B=\C_+\vee\C_-$, and let
$t=|\sigma(d_Y)|$.

\begin{lemma}\label{lem:bal}
The class weights satisfy
\[
\left|\sum_\lambda u_\lambda\,\sigma(\theta_\lambda)\right|
=t\,\FP(Y).
\]
\end{lemma}

\begin{proof}
Balancing and Lemma~\ref{lem:galcol} give
\[
\sum_cN_{YY}^c\,d_c\theta_c
=\theta_Y^2\Hopf(Y,Y)
=\theta_Y^2d_Y\tau(d_Y).
\]
Apply $\sigma$ and take absolute values. Since twists are roots of unity,
$|\sigma(\theta_Y)|=1$. Also
$\sigma(\tau(d_Y))=d_Y=\FP(Y)$, so the right-hand side has absolute value
$t\,\FP(Y)$.
\end{proof}

\begin{proposition}\label{prop:align}
Every simple constituent of $Y\otimes Y$ has twist $1$ or $-1$.
\end{proposition}

\begin{proof}
By Lemma~\ref{lem:Yloc}, $Y$ belongs to the subcategory $\mathcal B$, which
has property F. Theorem~\ref{thm:ms} and Lemma~\ref{lem:bal} therefore give
\[
t\,\FP(Y)
=\left|\sum_\lambda u_\lambda\sigma(\theta_\lambda)\right|
\le\sum_\lambda|u_\lambda|
=V\le t\,\FP(Y).
\]
Equality holds throughout. Each weight is nonzero: indeed,
\[
u_\lambda=\sigma(m_\lambda),\qquad
m_\lambda:=\sum_{c\in\lambda}N_{YY}^c\,d_c>0,
\]
where positivity follows from pseudo-unitarity. Equality in the triangle
inequality consequently implies that there is $\omega\in\mathbb C$, with
$|\omega|=1$, such that
\[
\sigma(\theta_\lambda)=\operatorname{sgn}(u_\lambda)\,\omega
\qquad\text{for every class }\lambda.
\]
Since $Y$ is self-dual, the unit occurs in $Y\otimes Y$. Its twist is $1$,
so $\omega\in\{1,-1\}$. Hence every $\sigma(\theta_\lambda)$ belongs to
$\{1,-1\}$, and applying $\sigma^{-1}$ proves the assertion.
\end{proof}

\section{A Markov identity and a Galois obstruction}\label{sec:QO}

We first establish two trace identities in an arbitrary ribbon fusion
category. Let $A$ be simple and write $c=c_{A,A}$. For a simple constituent
$X$ of $A\otimes A$, the braiding acts on the multiplicity space by
postcomposition,
\[
c_X:\Hom(X,A\otimes A)\to\Hom(X,A\otimes A),\qquad c_X(f)=c\circ f,
\]
and we set $t_k(X):=\tr(c_X^k)$. Balancing implies that $c_X^2$ is the scalar
$\theta_X\theta_A^{-2}$.

\begin{lemma}[Kink identity]\label{lem:kink}
One has
\[
\sum_Xd_X\,t_1(X)=\Tr_{A\otimes A}(c)=\theta_A d_A.
\]
\end{lemma}

\begin{proof}
The first equality follows by decomposing $A\otimes A$ into simple objects
with multiplicity spaces. For the second, take the right partial trace:
$\Tr(c)=\Tr_A(\ptr_2(c))=\theta_A d_A$.
\end{proof}

\begin{lemma}[Markov identity]\label{lem:markov}
One has
\[
\sum_Xt_1(X)\,\Hopf(X,A)
=\theta_A\Tr_{A\otimes A}(c^4)
=\theta_A\sum_Xd_X\,t_4(X).
\]
\end{lemma}

In a modular category, Lemma~\ref{lem:markov} is a consequence of the
self-braiding trace formula of Bantay \cite{Ban}, which expresses each $t_1(X)$
in terms of modular data: pairing that formula with the $S$-matrix column of
$A$ gives the lemma. See Giorgetti--Rehren \cite[Proposition~2.3 of the arXiv version]{GR} for a
proof in unitary modular categories, and Ng--Rowell--Wen \cite[Theorem~3]{NRW}
for the premodular setting and its relation to generalised Frobenius--Schur
indicators. The lemma itself holds in any ribbon fusion category, and we
include a short braid-group proof. We compute the trace of one three-strand braid
in two ways. Decomposing the
first two strands into simple constituents produces the Hopf-link sum on the
left. A braid conjugation followed by a partial trace produces
$\theta_A\Tr(c^4)$ on the right.

\begin{proof}
Let $a=\sigma_1$ and $b=\sigma_2$ act on $A^{\otimes3}$, and put
$\mu=ba^2b$. By the hexagon identities, $\mu$ is the block monodromy
$c_{A,A\otimes A}\circ c_{A\otimes A,A}$: the first two strands, regarded as a
single block, travel once around the third strand. On the summand
corresponding to $\Hom(X,A\otimes A)\otimes(X\otimes A)$, naturality
identifies the actions of $a$ and $\mu$ with
\[
c_X\otimes\mathrm{id}_{X\otimes A},
\qquad
\mathrm{id}\otimes(c_{A,X}c_{X,A}),
\]
respectively. Taking categorical traces gives
\begin{equation}\label{eq:block-trace}
\Tr(a\mu)=\sum_Xt_1(X)\,\Hopf(X,A).
\end{equation}

The braid relation gives an explicit conjugacy. For $w=a\mu=aba^2b$,
\[
w=babab,\qquad
b^{-1}wb=abab^2=bab^3,\qquad
b^{-2}wb^2=ab^4.
\]
Conjugation by $\Delta=aba$ exchanges $a$ and $b$, so $w$ is conjugate to
$ba^4$. Categorical trace is invariant under conjugacy. Taking the partial
trace over the third strand therefore yields
\[
\begin{aligned}
\Tr(a\mu)&=\Tr(ba^4)\\
&=\Tr_{A\otimes A}\bigl((\mathrm{id}_A\otimes\ptr_2(c))c^4\bigr)\\
&=\theta_A\Tr_{A\otimes A}(c^4).
\end{aligned}
\]
Together with \eqref{eq:block-trace} and the decomposition into multiplicity
spaces, this proves the identity.
\end{proof}

\begin{theorem}\label{thm:QO}
Let $\M$ be a modular category, $\tau$ a Galois automorphism, and
$A:=\htau(\one)$. If every simple constituent of $A\otimes A$ has twist $1$ or
$-1$, then $d_A^2=1$.
\end{theorem}

Recall that $A$ is self-dual by Proposition~\ref{prop:galois}, so
$A\otimes A=A\otimes A^*$. The idea of the proof is as follows. The twist
assumption restricts the eigenvalues of the rescaled braiding to fourth roots of
unity. The kink and Markov identities then give two expressions with the same
absolute value, one involving the dimensions $d_X$ and the other their Galois
conjugates. Comparing them forces $d_A^2$ to be fixed by $\tau$, and the Galois
identity~\eqref{eq:galois-global} then forces $d_A^2=1$. Pseudo-unitarity is not
used.

\begin{proof}
On the multiplicity space $\Hom(X,A\otimes A)$, balancing gives
\[
(\theta_Ac_X)^2=\theta_X\,\mathrm{id}\in\{\mathrm{id},-\mathrm{id}\}.
\]
Thus $\theta_Ac_X$ is diagonalisable with eigenvalues in $\{1,-1,i,-i\}$.
Let $n_\mu(X)$ be the multiplicity of the eigenvalue $\mu$. Then
\begin{equation}\label{eq:eigen-traces}
t_1(X)=\theta_A^{-1}\sum_\mu\mu\,n_\mu(X),\qquad
t_4(X)=\theta_A^{-4}N_{AA}^X.
\end{equation}
Define the real numbers
\[
r:=\sum_X\bigl(n_1(X)-n_{-1}(X)\bigr)d_X,\qquad
s:=\sum_X\bigl(n_i(X)-n_{-i}(X)\bigr)d_X.
\]
Substituting \eqref{eq:eigen-traces} into the kink identity gives the first
equation below. Substituting \eqref{eq:eigen-traces} and the Galois-column
identity $\Hopf(X,A)=d_A\tau(d_X)$ of Lemma~\ref{lem:galcol} into the Markov
identity gives
$\theta_A^{-1}d_A\bigl(\tau(r)+i\tau(s)\bigr)=\theta_A^{-3}\sum_XN_{AA}^Xd_X
=\theta_A^{-3}d_A^2$, which is the second:
\begin{equation}\label{eq:parallel}
\begin{aligned}
r+is&=\theta_A^2\,d_A&&\text{(kink identity)},\\
\tau(r)+i\tau(s)&=\theta_A^{-2}\,d_A&&\text{(Markov identity)}.
\end{aligned}
\end{equation}
In the Markov identity the dimension weights $d_X$ are replaced by
$\tau(d_X)$, while the eigenvalue labels $1,-1,i,-i$ remain fixed. This is why
the second equation contains $i\tau(s)$, even when $\tau(i)=-i$.

Dimensions in a spherical fusion category are real, and so are their Galois
conjugates by~\eqref{eq:galois-dimension}. Since $r$, $s$, $\tau(r)$ and
$\tau(s)$ are therefore real and $|\theta_A|=1$, taking
squared absolute values in \eqref{eq:parallel} gives
\[
r^2+s^2=d_A^2,\qquad \tau(r)^2+\tau(s)^2=d_A^2.
\]
Applying $\tau$ to the first equality shows that
\begin{equation}\label{eq:fixed-square}
\tau(d_A^2)=d_A^2.
\end{equation}

Write $D=\dim(\M)$ and $\rho=d_A^2=D/\tau(D)$, by~\eqref{eq:galois-global}.
Choose $m\ge1$ such that $\tau^m$ fixes the cyclotomic field $K$. Since
$\tau(\rho)=\rho$ by~\eqref{eq:fixed-square}, the product telescopes:
\[
\rho^m=\prod_{j=0}^{m-1}\tau^j(\rho)
=\prod_{j=0}^{m-1}\frac{\tau^j(D)}{\tau^{j+1}(D)}
=\frac{D}{\tau^m(D)}=1.
\]
Since $d_A$ is real and nonzero, $\rho>0$, so $\rho=1$.
\end{proof}

\begin{corollary}\label{cor:QO}
If, in Theorem~\ref{thm:QO}, $\M$ is pseudo-unitary, then $A$ is invertible.
Equivalently, in a pseudo-unitary modular category, every non-invertible
image $\htau(\one)$ of the unit label has a simple constituent of
$\htau(\one)\otimes\htau(\one)$ whose twist is not $\pm1$.
\end{corollary}

\begin{proof}
Pseudo-unitarity gives $\FP(A)=d_A=1$, so $A$ is invertible.
\end{proof}

\begin{remark}\label{rem:hypotheses}
The twist hypothesis makes both traces in \eqref{eq:eigen-traces} expressible
through integer eigenvalue multiplicities. The Galois-column identity then
relates the kink and Markov identities. Pseudo-unitarity enters only in
Corollary~\ref{cor:QO}, and that step needs genuine extra input: outside the
pseudo-unitary setting, categorical dimension $\pm1$ does not imply
invertibility. For example, in $\mathrm{Fib}\boxtimes\mathrm{Fib}^\sigma$, with
$\sigma(\varphi)=-\varphi^{-1}$, the product $A$ of the two nontrivial simple
objects lies in the Galois orbit of the unit and has $d_A=-1$ but
$\FP(A)=\varphi^2$ \cite{PSYZ}. (This $A$ does not satisfy the twist hypothesis,
so it is not a counterexample to a stronger statement.)

The Galois hypothesis on $A$ cannot be omitted
from the theorem. For example, let $SU(4)_4$ be the modular category
associated with the Lie algebra of $SU(4)$ at level $4$, and let
$\mathrm{adj}$ denote its adjoint object. The fundamental object $V$
satisfies the fundamental--dual fusion rule
\[
V\otimes V^*=\one\oplus\mathrm{adj}.
\]
The dimensions and twists in \sourcecite{DGPRZ}{Table~2}{https://arxiv.org/html/2005.05544v2\#S6.T2} give
$\theta_{\mathrm{adj}}=-1$ and $d_V^2=4+2\sqrt2$.
Thus all simple constituents of $V\otimes V^*$ have twist $\pm1$, although
$V$ is non-invertible. The object $V$ is not self-dual.

Nor can $\{1,-1\}$ be replaced by the fourth roots of unity. In $SU(2)_6$ the
spin-one object $A=X_2$ has $d_A=1+\sqrt2$, and by the truncated
Clebsch--Gordan rule $A\otimes A=\one\oplus X_2\oplus X_4$, with twists $1$, $i$
and $-i$ \cite[\S4]{BRWZ}. Yet $A$ is a Galois image of the unit, obtained from
$\zeta_{16}\mapsto\zeta_{16}^3$ by the sine formula for the $S$-matrix.
\end{remark}

\section{Proof of the main result}\label{sec:main}

\begin{proof}[Proof of Theorem~\ref{thm:main}]
Suppose that $\C$ has a pseudo-unitary Galois conjugate and property F, but
is not weakly integral. By Lemma~\ref{lem:galois-transfer} we may replace $\C$
by that conjugate, so $\C$ is pseudo-unitary. Proposition~\ref{prop:red} shows
that the modular category $\M=\mathcal Z(\C)$ is pseudo-unitary and not weakly
integral, and that $\mathcal B=\C_+\vee\C_-$ has property F. The construction in
Section~\ref{sec:red} gives a non-invertible self-dual object $Y=\htau(\one)$
in $\mathcal B$. By Proposition~\ref{prop:align}, every simple constituent of
$Y\otimes Y$ has twist $1$ or $-1$. Then Corollary~\ref{cor:QO} shows that $Y$
is invertible, a contradiction.
\end{proof}

The theorem also identifies infinite braid images object by object. This
answers the question someone studying a particular object will ask, and it
addresses the objectwise form of the conjecture
\cite[Conjecture~1.5]{GHR}.

\begin{corollary}\label{cor:object}
Let $X$ be a simple object of a braided fusion category, and suppose that the
fusion subcategory $\langle X\rangle$ has a pseudo-unitary Galois conjugate. If
$\FP(X)^2\notin\Z$, then $G_{X,n}$ is infinite for every sufficiently large $n$.
In particular, this applies to every such simple object of a pseudo-unitary
braided fusion category.
\end{corollary}

\begin{proof}
If every $G_{X,n}$ were finite, then $\langle X\rangle$ would have property F by
Lemma~\ref{lem:closure}, and so would be weakly integral by
Theorem~\ref{thm:main}. In a weakly integral fusion category every simple object
has integer squared Frobenius--Perron dimension \cite[Proposition~8.27]{ENO}, a
contradiction. Once
$G_{X,n_0}$ is infinite, tensoring with identities embeds it in $G_{X,n}$ for
every $n\ge n_0$.
\end{proof}

The word \emph{simple} matters: a direct sum of simple objects in a weakly
integral category can have non-integral squared Frobenius--Perron dimension.

The proof is also effective. In the setting of Section~\ref{sec:align}, with
$\M$ pseudo-unitary and $Y=\htau(\one)$ non-invertible, Corollary~\ref{cor:QO}
gives a simple constituent of $Y\otimes Y$ with twist outside $\{1,-1\}$, while
the unit is a constituent with twist $1$. Since every class weight $u_\lambda$
is nonzero, the triangle inequality in Lemma~\ref{lem:bal} is strict:
$V>t\,\FP(Y)$. Corollary~\ref{cor:strands} therefore produces an explicit strand
number $n_0$, computed from the modular data alone, such that $G_{Y,n}$ is
infinite for every $n\ge n_0$. In computations one can certify a smaller $n$ by
checking $V^{n-1}>t^{n-2}M_n$ directly, with $M_n$ computed from powers of the
fusion matrix.

\section{Consequences and remaining questions}\label{sec:conseq}

\subsection{Subcategories}
Property F passes to fusion subcategories, so the theorem also gives the
following obstruction without any Galois-conjugacy hypothesis on the ambient
category.

\begin{corollary}
Let $\C$ be a braided fusion category containing a fusion subcategory $\mathcal A$
with a pseudo-unitary Galois conjugate. If $\mathcal A$ is not weakly integral,
then $\C$ does not have property F.
\end{corollary}

\begin{proof}
The braiding of $\C$ restricts to $\mathcal A$. If $\C$ had property F, then
so would $\mathcal A$, contradicting Theorem~\ref{thm:main}.
\end{proof}

Similarly, if $\C$ is generated by fusion subcategories $\mathcal A_i$, each with
a pseudo-unitary Galois conjugate (possibly for different automorphisms), and
$\C$ has property F, then every $\mathcal A_i$ is weakly integral. So is $\C$.
Indeed, for fusion subcategories $\mathcal A,\mathcal B$ of a braided fusion
category, $\FP(\mathcal A)\FP(\mathcal B)=\FP(\mathcal A\vee\mathcal B)\,
\FP(\mathcal A\cap\mathcal B)$ \cite[Lemma~3.38]{DGNO}. If $\mathcal A$ and
$\mathcal B$ are weakly integral, so is $\mathcal A\cap\mathcal B$, as its simple
objects are simple in $\mathcal A$; hence $\FP(\mathcal A\vee\mathcal B)$ is a
rational algebraic integer, so an integer. Induction on the number of
subcategories finishes the argument.

\subsection{Beyond the pseudo-unitary hypothesis}
Theorem~\ref{thm:ms} needs neither pseudo-unitarity nor self-duality, so it
gives infinite-image certificates outside the hypothesis of
Theorem~\ref{thm:main}. Earlier trace and positivity obstructions of a related
kind appear in Rowell's work on non-unitarisable ribbon categories
\cite{Row04}. For example, let
$\mathcal R=\mathcal C(\mathfrak{so}_5,3/2)_{\mathrm{ad}}$, the rank-six
adjoint subcategory at $q=e^{\pi i/9}$, whose fusion ring has no pseudo-unitary
categorification \cite[Theorem~1]{Sch}. Label its simple objects
$X_0,\ldots,X_5$ as in \cite{Sch,Row05}, and set $a=2\cos(\pi/9)$ and $t=2\cos(4\pi/9)$. Then $Y=X_2$ satisfies
\[
Y\otimes Y\cong\one\oplus X_3\oplus X_5,\qquad
Y^{\otimes3}\cong3X_2\oplus X_3\oplus2X_4,
\]
with categorical dimensions $d_{X_2}=t$, $d_{X_3}=1$, $d_{X_4}=-1$ and
$d_{X_5}=-a$, and the three constituents of $Y\otimes Y$ have distinct twists.
Taking $\sigma=\mathrm{id}$ gives $V=2+a$ and $M_3=3t+3$. Using $a=2-t^2$ and
$t^3=3t-1$,
\[
V^2-tM_3=16-8t^2-4t>0,
\]
since $0<t<\tfrac12$. By Theorem~\ref{thm:ms} with $n=3$, the image $G_{Y,3}$
is infinite. This is not a new infinite-image result; it illustrates the reach
of the trace obstruction.

Removing pseudo-unitarity from Theorem~\ref{thm:main} needs a new idea. Without
it, Lemma~\ref{lem:bal} only gives
$\bigl|\sum_\lambda u_\lambda\sigma(\theta_\lambda)\bigr|=t\,|d_Y|$, whereas
Corollary~\ref{cor:strands} gives $V\le t\,\FP(Y)$. When $|d_Y|<\FP(Y)$ the
bounds no longer meet, and a twist class with zero total dimension is invisible
to the trace calculation.

\subsection{Further trace inequalities}
The proof of Theorem~\ref{thm:ms} extends in two directions. For simple $X$ and
$Y$, group the constituents of $X\otimes Y$ by twist and let $V_\sigma(X,Y)$ be
the sum of the absolute values of the class traces in $\mathcal B^\sigma$. Using
alternating strands $X,Y,X,Y,\ldots$ and the pure braid group, property F gives
\[
V_\sigma(X,Y)^2\le|\sigma(d_X)\sigma(d_Y)|\,\FP(X)\FP(Y).
\]
Similarly, using spectral projections of the full twist on $Y^{\otimes k}$ and
blocks of $k$ strands that share one strand with their neighbours, property F
gives $V_k\le t\,\FP(Y)^{k-1}$, where $V_k$ is the analogous sum for the
constituents of $Y^{\otimes k}$. Both inequalities use only fusion rules,
dimensions and twists.

\subsection{The converse direction}
These methods cannot prove the converse. If $\C$ is weakly integral, then
$|\sigma(d_W)|=d_W$ for all simple $W$, every class weight is positive, and
$V=d_Y^2$, $t=d_Y$ and $M_n=d_Y^n$. So Theorem~\ref{thm:ms} holds with
equality, independently of the braid image. Likewise every Galois image of the
unit is invertible, since $\dim(\C)\in\Z$. The implication from weak integrality
to property F is known for weakly group-theoretical categories
\cite[Theorem~7.2]{GN}. For weakly group-theoretical modular categories with
their canonical positive spherical structure, Li and Ruan have shown that the
Reshetikhin--Turaev mapping class group representations have finite projective
image \cite{LR}. In general the implication would follow from a negative answer
to \cite[Question~2]{ENO2}. The argument here also
leaves open the implication from property F to weak integrality for braided
fusion categories with no pseudo-unitary Galois conjugate.

\begin{thebibliography}{DGPRZ}

\bibitem[Ban]{Ban}
P. Bantay,
\emph{The Frobenius--Schur indicator in conformal field theory},
Phys. Lett. B \textbf{394} (1997), 87--88.

\bibitem[BRWZ]{BRWZ}
P. Bonderson, E. C. Rowell, Z. Wang and Q. Zhang,
\href{https://arxiv.org/abs/1704.02041}{\emph{Congruence subgroups and super-modular categories}},
arXiv:1704.02041.

\bibitem[CG]{CG}
A. Coste and T. Gannon,
\href{https://doi.org/10.1016/0370-2693(94)91226-2}{\emph{Remarks on Galois symmetry in rational conformal field theories}},
Phys. Lett. B \textbf{323} (1994), 316--321.

\bibitem[Del]{Del}
P. Deligne,
\emph{Cat\'egories tensorielles},
Mosc. Math. J. \textbf{2} (2002), 227--248.

\bibitem[DGNO]{DGNO}
V. Drinfeld, S. Gelaki, D. Nikshych and V. Ostrik,
\href{https://arxiv.org/abs/0906.0620}{\emph{On braided fusion categories I}},
Selecta Math. (N.S.) \textbf{16} (2010), 1--119.

\bibitem[DGPRZ]{DGPRZ}
C. Delaney, C. Galindo, J. Plavnik, E. C. Rowell and Q. Zhang,
\href{https://arxiv.org/abs/2005.05544}{\emph{Braided zesting and its applications}},
Comm. Math. Phys. \textbf{386} (2021), 1--55.

\bibitem[DLN]{DLN}
C. Dong, X. Lin and S.-H. Ng,
\href{https://arxiv.org/abs/1201.6644}{\emph{Congruence property in conformal field theory}},
Algebra Number Theory \textbf{9} (2015), 2121--2166.

\bibitem[E]{E}
P. Etingof,
\href{https://arxiv.org/abs/math/0207007}{\emph{On Vafa's theorem for tensor categories}},
Math. Res. Lett. \textbf{9} (2002), 651--657.

\bibitem[EGNO]{EGNO}
P. Etingof, S. Gelaki, D. Nikshych and V. Ostrik,
\href{https://doi.org/10.1090/surv/205}{\emph{Tensor Categories}},
Mathematical Surveys and Monographs, vol.~205,
American Mathematical Society, Providence, RI, 2015.

\bibitem[ENO]{ENO}
P. Etingof, D. Nikshych and V. Ostrik,
\href{https://doi.org/10.4007/annals.2005.162.581}{\emph{On fusion categories}},
Ann. of Math. (2) \textbf{162} (2005), 581--642.

\bibitem[ENO2]{ENO2}
P. Etingof, D. Nikshych and V. Ostrik,
\href{https://arxiv.org/abs/0809.3031}{\emph{Weakly group-theoretical and solvable fusion categories}},
Adv. Math. \textbf{226} (2011), 176--205.

\bibitem[ERW]{ERW}
P. Etingof, E. Rowell and S. Witherspoon,
\href{https://arxiv.org/abs/math/0703274}{\emph{Braid group representations from twisted quantum doubles of finite groups}},
Pacific J. Math. \textbf{234} (2008), 33--41.

\bibitem[FLW]{FLW}
M. H. Freedman, M. J. Larsen and Z. Wang,
\href{https://arxiv.org/abs/math/0103200}{\emph{The two-eigenvalue problem and density of Jones representation of braid groups}},
Comm. Math. Phys. \textbf{228} (2002), 177--199.

\bibitem[GHR]{GHR}
C. Galindo, S.-M. Hong and E. C. Rowell,
\href{https://arxiv.org/abs/1105.5048}{\emph{Generalized and quasi-localizations of braid group representations}},
arXiv:1105.5048.

\bibitem[GN]{GN}
J. Green and D. Nikshych,
\href{https://arxiv.org/abs/1911.02633}{\emph{On the braid group representations coming from weakly group-theoretical
fusion categories}},
J. Algebra Appl. \textbf{20} (2021), article no.~2150210.

\bibitem[GR]{GR}
L. Giorgetti and K.-H. Rehren,
\href{https://arxiv.org/abs/1606.04378}{\emph{Bantay's trace in unitary modular tensor categories}},
Adv. Math. \textbf{319} (2017), 211--223; arXiv:1606.04378.

\bibitem[LR]{LR}
E. Li and Y. Ruan,
\href{https://arxiv.org/abs/2610.05177}{\emph{Generalized property F for weakly group-theoretical modular fusion categories}},
arXiv:2610.05177.

\bibitem[LRW]{LRW}
M. J. Larsen, E. C. Rowell and Z. Wang,
\href{https://arxiv.org/abs/math/0506025}{\emph{The $N$-eigenvalue problem and two applications}},
Int. Math. Res. Not. \textbf{2005} (2005), 3987--4018.

\bibitem[NR]{NR}
D. Naidu and E. C. Rowell,
\href{https://arxiv.org/abs/0903.4157}{\emph{A finiteness property for braided fusion categories}},
Algebr. Represent. Theory \textbf{14} (2011), 837--855.

\bibitem[NRW]{NRW}
S.-H. Ng, E. C. Rowell and X.-G. Wen,
\href{https://arxiv.org/abs/2408.02748}{\emph{Recovering $R$-symbols from modular data}},
arXiv:2408.02748.

\bibitem[NS]{NS}
S.-H. Ng and P. Schauenburg,
\href{https://arxiv.org/abs/0806.2493}{\emph{Congruence subgroups and generalized Frobenius--Schur indicators}},
Comm. Math. Phys. \textbf{300} (2010), 1--46.

\bibitem[PSYZ]{PSYZ}
J. Plavnik, A. Schopieray, Z. Yu and Q. Zhang,
\href{https://arxiv.org/abs/2111.05228}{\emph{Modular tensor categories, subcategories, and Galois orbits}},
arXiv:2111.05228.

\bibitem[Row04]{Row04}
E. C. Rowell,
\href{https://arxiv.org/abs/math/0403217}{\emph{On a family of non-unitarizable ribbon categories}},
arXiv:math/0403217.

\bibitem[Row05]{Row05}
E. C. Rowell,
\href{https://arxiv.org/abs/math/0503226}{\emph{From quantum groups to unitary modular tensor categories}},
in \emph{Representations of algebraic groups, quantum groups, and Lie algebras},
Contemp. Math. \textbf{413}, Amer. Math. Soc., 2006, 215--230.

\bibitem[Sch]{Sch}
A. Schopieray,
\href{https://arxiv.org/abs/2010.02958}{\emph{Non-pseudounitary fusion}},
arXiv:2010.02958.

\end{thebibliography}

\end{document}
